\subsection{General quadratic variation and exact \texorpdfstring{$j^1$}{j1}}
For bounded-jump martingales we have finite-horizon $M^2$ quadratic processes,
stopping consistency, gluing, uniqueness, terminal roots, and the one-sided
Davis estimate with constant $6$. Horizonwise jump bounds suffice for local
martingales; the DDY decomposition supplies them for its component $L^T$.

\begin{lemma}[Absolute convergence of finite-variation jump corrections]
For the same DDY decomposition $N=L+Q$, let $V_T=\Var_{[0,T]}Q$.
On one full-measure set, for every $T\geq0$,
\[
 \sum_{0<t\leq T}|\Delta Q_t|\leq V_T,\quad
 \sum_{0<t\leq T}|\Delta L_t\Delta Q_t|\leq2cV_T,\quad
 \sum_{0<t\leq T}(\Delta Q_t)^2\leq V_T^2.
\]
Moreover the absolutely convergent correction is
\[
 D_T=\sum_{0<t\leq T}\bigl((\Delta N_t)^2-(\Delta L_t)^2\bigr)
 =2\sum_{0<t\leq T}\Delta L_t\Delta Q_t+\sum_{0<t\leq T}(\Delta Q_t)^2,
 \qquad |D_T|\leq4cV_T+V_T^2.
\]
\end{lemma}
\begin{proof}
For a right-continuous BV path, the variation of its signed Stieltjes measure
at a singleton is the absolute jump. Every finite sum of absolute jumps in
$(0,T]$ is bounded by the measure's variation there, proving summability and
the first estimate. Apply this to a locally BV path stopped at $T$.
Multiply by $|\Delta L|\leq2c$ and $|\Delta Q|\leq V_T$ for the other bounds.
Expand $\Delta N=\Delta L+\Delta Q$ and add the absolutely convergent series.
The common all-time jump-bound event works for every $T$.
\end{proof}
\begin{lemma}[Squared increments on the same finite grid]
For a finite ordered grid $u_0\leq\cdots\leq u_m$, define
\[
 \begin{aligned}
 E_G(Y,t)&=\sum_{k<m}(Y_{t\wedge u_{k+1}}-Y_{t\wedge u_k})^2,\\
 C_G(Y,Z,t)&=\sum_{k<m}(Y_{t\wedge u_{k+1}}-Y_{t\wedge u_k})
                       (Z_{t\wedge u_{k+1}}-Z_{t\wedge u_k}).
 \end{aligned}
\]
For the same DDY components, pathwise at every time,
\[
 E_G(N,t)=E_G(L,t)+2C_G(L,Q,t)+E_G(Q,t),
\]
\[
 |\sqrt{E_G(N,t)}-\sqrt{E_G(L,t)}|\leq\Var_{[0,t]}(Q).
\]
The estimate contains no grid-dependent constant.
\end{lemma}
\begin{proof}
Expanding the square of each cell increment gives the first identity.
Cauchy--Schwarz gives
$C_G(Y,Z,t)\leq\sqrt{E_G(Y,t)}\sqrt{E_G(Z,t)}$.
Together with the square expansion this proves the root triangle inequality
$\sqrt{E_G(Y+Z,t)}\leq\sqrt{E_G(Y,t)}+\sqrt{E_G(Z,t)}$.
The sum of absolute increments over ordered cells is at most path variation,
so
\[
 \sqrt{E_G(Q,t)}\leq\sum_{k<m}|Q_{t\wedge u_{k+1}}-Q_{t\wedge u_k}|
 \leq\Var_{[0,t]}(Q).
\]
Apply the root triangle inequality to both $N=L+Q$ and $L=N-Q$.
\end{proof}
\begin{lemma}[Factorial cross-increment limits]
Let $f$ be a real c\`adl\`ag path and $Q$ a right-continuous locally BV path.
Let $\eta_T$ be the signed Stieltjes measure of $Q^T_t=Q_{t\wedge T}$ and
$G_r$ the grid $k/(r+1)!$, $0\leq k\leq(r+1)(r+1)!$. Then
\[
 C_{G_r}(f,Q,T)\longrightarrow\int_{(0,T]}\Delta f_s\,d\eta_T(s).
\]
Consequently, for the same DDY decomposition, at every path and finite horizon,
\[
 E_{G_r}(N,T)-E_{G_r}(L,T)\longrightarrow
 2\int_{(0,T]}\Delta L_s\,d\eta_T(s)+\int_{(0,T]}\Delta Q_s\,d\eta_T(s).
\]
\end{lemma}
\begin{proof}
For the half-open cell $(a_r(s),b_r(s)]$ containing $s>0$, put
$D_rf(s)=f(b_r(s)\wedge T)-f(a_r(s)\wedge T)$.
The endpoints approach $s$ with $a_r(s)<s\leq b_r(s)$, so
$D_rf(s)\to\Delta f_s$ for $s\leq T$. This also holds if $s$ itself is a
mesh point. Compact boundedness of c\`adl\`ag paths gives
$|D_rf(s)|\leq2\sup_{u\leq T}|f(u)|$.
Dominated convergence under the finite measure $|\eta_T|$ makes the
$L^1$ error over $(0,T]$ tend to zero, hence the signed integrals converge.
The integral on each stopped cell equals
$(f(b\wedge T)-f(a\wedge T))(Q(b\wedge T)-Q(a\wedge T))$.
For sufficiently large $r$, the cells cover $(0,T]$, so the integral is
the cross-increment sum. Apply this with $f=L$ and $f=Q$ and use the finite
square expansion to obtain the second limit.
\end{proof}
\begin{lemma}[Jump integrals as series]
For a real c\`adl\`ag path $f$ and a finite signed measure $\eta$,
\[
 \int_{(0,T]}\Delta f_t\,d\eta(t)
 =\sum_{t\in(0,T]}\Delta f_t\,\eta(\{t\}).
\]
In particular, for $\eta_T$ and the same DDY components, pathwise,
\[
 E_{G_r}(N,T)-E_{G_r}(L,T)\longrightarrow
 2\sum_{0<t\leq T}\Delta L_t\Delta Q_t+
 \sum_{0<t\leq T}(\Delta Q_t)^2.
\]
\end{lemma}
\begin{proof}
For each $n$, only finitely many times in $(0,T]$ have
$|\Delta f_t|>1/(n+1)$. Thus the support of nonzero jumps is countable.
The uniform bound on the measurable cell approximations passes to their
limit $\Delta f$, making it bounded on $(0,T]$ and integrable for $|\eta|$.
Decompose the support into disjoint singletons and sum the integrals to
obtain the first identity. For a stopped FV path,
$\eta_T(\{t\})=\Delta Q_t$ whenever $0<t\leq T$, including the terminal
jump. Apply the identity with $f=L,Q$ and rewrite the preceding integral
limit to obtain the series formula.
\end{proof}
\begin{lemma}[Nonnegativity of the corrected candidate]
Let $V$ now denote the glued quadratic variation of $L$. On one full-measure
set, $V_t+D_t\geq0$ for every $t$.
\end{lemma}
\begin{proof}
On coordinate $n$, at times before its localizer and horizon, the original
and stopped squared-increment sums agree on the same factorial grids.
Use exactly the forward weights and cutoff retained by the construction
of the quadratic variation of $L^{\tau_n}$. Full-sequence convergence of the
correction is preserved by these weights. Hence the convex squared sums of
$N^{\tau_n}$ tend to the coordinate variation plus $D_t$ and remain
nonnegative. Coordinate agreement with $V$ and exhaustion remove the stops.
Take the countable intersection of coordinate events before choosing $t$.
\end{proof}
\begin{lemma}[Regularity of the correction]
The processes $D$ and $V+D$ are strongly adapted and are right-continuous
and locally of finite variation on one full-measure set.
\end{lemma}
\begin{proof}
On a finite horizon, the coefficients
$a_s=2\Delta L_s\Delta Q_s+(\Delta Q_s)^2$ are absolutely summable.
Dominated convergence for the series makes the cumulative sum
$\sum_{s\leq t}a_s$ right-continuous. Write
$a_s=(a_s+|a_s|)-|a_s|$; both cumulative sums are bounded increasing functions,
so their difference has finite variation. These cumulative sums locally
agree with $D$. Adaptation follows from the pathwise limit of the measurable
finite-grid square differences at each time. Add the regular process $V$.
\end{proof}
\begin{lemma}[Monotonicity of the candidate]
On one full-measure set $V+D$ is increasing at all times.
\end{lemma}
\begin{proof}
Use the same convex convergence as in the nonnegativity argument. For $s\leq t$
choose a coordinate whose localizer exceeds $t+1$. Approximate $s,t$ from the
right by capped factorial grid points $u_r\leq v_r$ below that localizer.
The tail condition on the weights ensures that later convex sums use grids
refining this common coarse grid. Squared sums are increasing between grid
points, so convex averaging preserves their order. Pass first to the convex
limit, identify with the candidate, and then let $r\to\infty$ by right
continuity. Coordinate events were intersected before choosing $s,t$.
\end{proof}
\begin{lemma}[An everywhere regular representative]
Under the usual conditions there is a strongly adapted $A\sim V+D$ whose
paths are right-continuous, increasing, locally BV, have left limits, and
satisfy $A_0=0$.
\end{lemma}
\begin{proof}
The set where right continuity or monotonicity fails is null, hence measurable
at time zero. Replace entire paths there by zero. This preserves adaptation
and indistinguishability. Initial values are zero by the empty jump sum and
zero start of $V$. Increasing paths have local bounded variation and left limits.
\end{proof}
\begin{lemma}[Jump-square identity]
On one full-measure set, $\Delta A_t=(\Delta N_t)^2$ for all times.
\end{lemma}
\begin{proof}
For an absolutely summable local family $a_s$, dominated convergence as
$u\uparrow t$ shows that the left limit of $\sum_{0<s\leq u}a_s$ is
$\sum_{0<s<t}a_s$. Thus its jump at $t>0$ is $a_t$.
Apply this to the correction and use $\Delta V_t=(\Delta L_t)^2$:
\[
 \Delta(V+D)_t=(\Delta L_t)^2+2\Delta L_t\Delta Q_t+(\Delta Q_t)^2
 =(\Delta N_t)^2.
\]
Both sides are zero at time zero. Transfer the all-time identity to $A$
by indistinguishability.
\end{proof}
\begin{lemma}[Closed stopping of the correction]
For a finite cutoff $\tau$, the correction formed from $L^\tau,Q^\tau$ equals
$D^\tau$. If $V^n$ is the quadratic coordinate at localizer $\tau_n$, then
\[
 A^{\tau_n}\sim(V^n)^{\tau_n}
 +2\sum_{0<s\leq\cdot}\Delta L^{\tau_n}_s\Delta Q^{\tau_n}_s
 +\sum_{0<s\leq\cdot}(\Delta Q^{\tau_n}_s)^2.
\]
\end{lemma}
\begin{proof}
Closed-stop jumps equal original jumps for $s\leq\tau$ and are zero afterward.
Summing to $t$ is therefore summing original jumps to $t\wedge\tau$, including
the boundary. Combine this pathwise identity with coordinate agreement for
$V$ and $A\sim V+D$.
\end{proof}
\begin{lemma}[The finite-variation square residual]
For a c\`adl\`ag locally BV path $Q$ and finite $T$,
\[
 Q_T^2-Q_0^2-\sum_{0<s\leq T}(\Delta Q_s)^2
 =2\int_{(0,T]}Q_{s-}\,dQ_s.
\]
\end{lemma}
\begin{proof}
Factorial left-step approximations converge to $Q_{s-}$, since the left cell
endpoint approaches every positive time strictly from below.
Compact boundedness and finite Stieltjes variation give dominated convergence.
On each finite grid the identity

\[
 2\sum Q_u(Q_v-Q_u)+\sum(Q_v-Q_u)^2=Q_T^2-Q_0^2
\]
 telescopes, with endpoints
capped at $T$. The preceding cross-increment and atomic integral results
identify the square-sum limit. Pass to the limit.
\end{proof}
For the same DDY components, put $I_Q(t)=\int_{(0,t]}Q_{s-}\,dQ_s$.
Since $Q_0=0$, on a common full-measure set,
\[
 N_t^2-A_t=(L_t^2-V_t)+
 2\left(L_tQ_t-\sum_{0<s\leq t}\Delta L_s\Delta Q_s+I_Q(t)\right).
\]
We next localize and pass martingale identities through dominated limits.
\begin{lemma}[Common localization with integrable variation]
For the same $L,Q$, choose finite $\rho_n\uparrow\infty$, $\rho_n\leq n+1$,
such that $L^{\rho_n},Q^{\rho_n}$ are true martingales,

\[
 |L^{\rho_n}_t|\leq n+1+2c,\qquad |Q_s|\leq n+1\quad(s<\rho_n).
\]
The first bound holds at all times on one full-measure set. Moreover,
$\Var_{[0,\rho_n]}Q$ is integrable and dominates every horizon variation of
$Q^{\rho_n}$.
\end{lemma}
\begin{proof}
The cumulative variation $A_t=\Var_{[0,t]}Q$ is adapted by finite-grid
representation, right-continuous and increasing. Take the minimum of true
localizers for both martingales, the level $n+1$ passages of $|L|$ and $A$,
and the horizon $n+1$. This is exhaustive and preserves true martingales,
using zero start to remove the time-zero indicators. The bounded jumps of
$L$ give the closed-stop bound. Before the stop $|Q_s|\leq A_s\leq n+1$.
With $a=n+1$, the variation jump identity gives
\[
 A_{\rho_n}\leq a+|\Delta Q_{\rho_n}|\leq2a+|Q_{\rho_n}|.
\]
The last stopped value is integrable by bounded optional sampling.
Measurability and this domination give integrability of $A_{\rho_n}$;
stopped variation equals $A_{t\wedge\rho_n}\leq A_{\rho_n}$.
\end{proof}
On this localization the finite-grid sums $I_G(Q^\rho,Q^\rho)$,
$I_G(Q^\rho,L^\rho)$, and $I_G(L^\rho,Q^\rho)$ are true martingales.
For the first two replace coefficients by $Q_s1_{\{s<\rho\}}$, which are
bounded and adapted; the sums agree because integrator increments after
stopping vanish. For the third use the closed-stop bound on $L$.
Neither boundedness nor square integrability of $Q_\rho$ is needed.

\begin{lemma}[Martingale limits under integrable domination]
If martingales $M^k$ converge almost surely at each time to strongly adapted
$Y$, and $|M^k_t|\leq B_t$ for one integrable $B_t$ at each time, then $Y$
is a martingale.
\end{lemma}
\begin{proof}
The limit is integrable by domination. For $s\leq t$ and $E\in\F_s$,
dominated convergence in the identities $\int_E M^k_s=\int_E M^k_t$ yields
$\int_EY_s=\int_EY_t$. Use the characterization of conditional expectation.
\end{proof}
\begin{lemma}[Self-integral and cross-term local martingales]
The processes
\[
 I_t=\int_{(0,t]}Q_{s-}\,dQ_s,\qquad
 K_t=L_tQ_t-\sum_{0<s\leq t}\Delta L_s\Delta Q_s
\]
are local martingales.
\end{lemma}
\begin{proof}
Fix a common localization and write $\bar L=L^{\rho_n}$, $\bar Q=Q^{\rho_n}$,
$a=n+1$, $b=n+1+2c$, and $A=\Var_{[0,\rho_n]}Q$.
Strict-prefix coefficients give the full self-integral bound
\[
 |I_G(\bar Q,\bar Q)_t|\leq a\Var_{[0,t]}(\bar Q)\leq aA.
\] Factorial sums converge to the stopped
Stieltjes integral. Compare the square identities before and after stopping
and use the stopped jump-square series to identify the limit with $I_{t\wedge\rho_n}$.
For cross terms, the finite-grid product formula with last grid point $u_G$ is
\[
 I_G(\bar L,\bar Q)_t+I_G(\bar Q,\bar L)_t
 =\bar L_{t\wedge u_G}\bar Q_{t\wedge u_G}-C_G(\bar L,\bar Q,t).
\]
The product is bounded by $bA$ and the cross sum by $2bA$, so the total is
bounded by $3bA$. Its factorial limit is $K_{t\wedge\rho_n}$ by the
cross-increment limit, the jump-series identity, and stopping compatibility.
The pathwise limits retain strong adaptation. Integrability of $A$ and the
preceding lemma make both closed stops true martingales. Zero start agrees
with the localizing indicator convention. We need not separately construct
a completed integral $\int Q_-\,dL$.
\end{proof}

\begin{theorem}[Quadratic variation of a local martingale]
Every zero-start strongly adapted c\`adl\`ag local martingale $N$ on a
filtered probability space satisfying the usual conditions has an adapted
c\`adl\`ag increasing process $V$ with
\[
 V_0=0,\qquad\Delta V=(\Delta N)^2,\qquad N^2-V\text{ a local martingale}.
\]
The jump identity holds simultaneously at all times almost surely, and $V$
is unique up to indistinguishability.
\end{theorem}
\begin{proof}
Use the regularized corrected candidate. The square formula gives
$N^2-V=(L^2-[L])+2(K+I)$.
Transfer the second summand to the indistinguishable regular process
$N^2-V-(L^2-[L])$, align localizers, and add the local martingales.
For uniqueness, allow even adapted c\`adl\`ag locally BV candidates $V,W$
without monotonicity, provided their initial values and jumps agree and both
square residuals are local martingales. Then $D=V-W$ is a locally BV local
martingale with zero jumps, hence has continuous paths almost surely.
Stop at integer horizons. The left-limit process is predictable and agrees
with this continuous process; correct the common null set and use predictable
FV local-martingale rigidity to make the difference zero. Countably many
horizons give all-time equality.
\end{proof}
Thus
\[
 [N]_T=[L]_T+D_T.
\] Different DDY thresholds and schedules give
indistinguishable certificates for the same source. Nonconvex full-sequence
square-sum convergence is not inferred from this existence result alone.
\begin{lemma}[Stopping consistency]
For any stopping time $\tau$, possibly infinite,
\[
 [N^\tau]\sim[N]^\tau.
\]
The left-hand side may be independently constructed.
\end{lemma}
\begin{proof}
A zero-start right-continuous local martingale remains one after arbitrary
closed stopping: stop each true localized martingale again, commute stops,
and remove the indicator at zero. Thus existence applies to $N^\tau$.
The candidate $V^\tau$ retains adaptation, c\`adl\`ag regularity,
monotonicity, and zero start. Before and at $\tau$ the jumps agree with the
original ones and afterward are zero. Also

\[
 (N^\tau)^2-V^\tau=(N^2-V)^\tau
\]
 is a local martingale.
Uniqueness for the stopped source proves the claim.
\end{proof}
\begin{lemma}[The root triangle inequality]
For zero-start strongly adapted c\`adl\`ag local martingales $X,Y$,

\[
 \sqrt{[X+Y]_t}\leq\sqrt{[X]_t}+\sqrt{[Y]_t}
\]
 at all times almost surely.
\end{lemma}
\begin{proof}
Write $V=[X]$, $W=[Y]$, $S=[X+Y]$. For real $a$ put

\[
 A^a=aS+(1-a)V+(a^2-a)W.
\] This is adapted c\`adl\`ag locally BV with zero
start, the correct squared jumps for $X+aY$, and square residual
\[
 (X+aY)^2-A^a=a((X+Y)^2-S)+(1-a)(X^2-V)+(a^2-a)(Y^2-W).
\]
The uniqueness comparison, which does not require monotonicity, gives
$A^a\sim[X+aY]$. Intersect the events for rational $a$.
At every time on this common event,
\[
 0\leq aS_t+(1-a)V_t+(a^2-a)W_t
 =W_ta^2+(S_t-V_t-W_t)a+V_t
\]
for all rational $a$, hence all real $a$ by continuity. Its discriminant gives
$(S_t-V_t-W_t)^2\leq4V_tW_t$. Since $V_t,W_t\geq0$, the root bound follows.
\end{proof}

\begin{lemma}[Finite-grid $L^1$ Davis estimates]
For a real discrete martingale $Y$, set

\[
 Y_n^*=\max_{0\leq k\leq n}|Y_k|,\qquad
 R_n=\sqrt{Y_0^2+\sum_{k<n}(Y_{k+1}-Y_k)^2}.
\]
Both are integrable, and
\[
 E Y_n^*\leq6E R_n,\qquad E R_n\leq7E Y_n^*.
\]
No square-integrability hypothesis is needed.
\end{lemma}
\begin{proof}
They are bounded respectively by $\sum_{k\leq n}|Y_k|$ and
$|Y_0|+\sum_{k<n}|Y_{k+1}-Y_k|$, proving integrability.
A bounded adapted coefficient times a martingale increment is integrable
with zero conditional expectation; hence the finite left-endpoint integral
has mean zero. Integrate the pathwise Davis inequality, whose coefficients
are bounded by one, to obtain the first estimate.
For the reverse inequality use, with $m\geq|x|$ and $m,q\geq0$,
\[
 F(x,m,q)=-2m+\sqrt{m^2+q}+\frac{m^2-x^2}{2\sqrt{m^2+q}},
\]
interpreting division by zero as zero. With $M=\max(m,|x+d|)$,
\[
 F(x+d,M,q+d^2)-F(x,m,q)
 \leq-\frac{xd}{\sqrt{m^2+q}}+5(M-m).
\]
To check this put $r=\sqrt{m^2+q}$ and $u=\sqrt{M^2+q+d^2}$.
For $r>0$ and $M\leq2m$, the square-root tangent estimate and $u\geq r$
bound the left side plus $xd/r$ by
$(M^2-m^2)/r-2(M-m)\leq M-m$.
For $M>2m$, $M=|x+d|$, $|d|\leq M+m\leq3(M-m)$, and $u-r\leq4(M-m)$.
Also $xd/r\leq|d|$. The new potential numerator $M^2-(x+d)^2$
is zero and the old numerator is nonnegative, giving the stated bound. If $r=0$ then $x=m=q=0$ and direct calculation suffices.
Sum with $x=Y_k,m=Y_k^*,q=R_k^2,d=Y_{k+1}-Y_k$.
Since the initial potential is nonpositive and the final one at least
$R_n-2Y_n^*$,
\[
 R_n\leq7Y_n^*-\sum_{k<n}
 \frac{Y_k}{\sqrt{R_k^2+(Y_k^*)^2}}(Y_{k+1}-Y_k).
\]
The coefficients are adapted and bounded by one; take expectations.
\end{proof}
Apply these estimates on the same DDY common localizers to the original
$N^{\rho_j}=L^{\rho_j}+Q^{\rho_j}$. Concavity prevents transferring the
reverse estimate directly to roots of convexified energies; we first need
convergence of the original square sums.
\begin{lemma}[Unconvexified square sums of a bounded martingale]
Under the usual conditions, let $M$ be a strongly adapted c\`adl\`ag
martingale satisfying $|M_t|\leq C$ on every path up to a finite horizon $T$.
On factorial grids $\pi_i$, put
\[
 I_i=2\sum_{[u,v]\in\pi_i}M_u(M_v-M_u),\qquad
 E_i=\sum_{[u,v]\in\pi_i}(M_v-M_u)^2.
\]
If $Y,V$ are the martingale and increasing components retained by the
quadratic construction, then $I_i\to Y_T$ and $E_i\to V_T$ in $L^2$.
If the bound holds almost surely at all times on one common set, then
$E_i\to V_T$ in probability as well.
\end{lemma}
\begin{proof}
Let $\nu$ be the predictable energy measure of the same $V$.
The grid left-endpoint coefficients $H_i$ are bounded predictable elementary
processes converging pointwise to $M_-$ on $(0,T]$. The measure $\nu$ is
finite and gives no mass to time zero or to times after $T$.
Dominated convergence yields $H_i\to M_-$ in $L^2(\nu)$.
The grid-independent isometry
\[
 \|((H_i-H_j)\cdot M)_T\|_{L^2(P)}
 =\|H_i-H_j\|_{L^2(\nu)}
\]
makes $I_i$ Cauchy in $L^2(P)$; denote its limit by $I$.
The retained forward convex weights preserve this limit. With those same
weights, however, the integral sums converge almost surely to $Y_T$.
Uniqueness of the probability limit gives $I=Y_T$.
Use $I_i+E_i=M_T^2-M_0^2$ and the retained square decomposition to obtain
$E_i\to V_T$.

If the bound fails only on one common null set, the usual conditions put
that set in $\F_0$. Replacing the process by zero there produces an
everywhere bounded indistinguishable c\`adl\`ag martingale.
The same weights and components supply quadratic data for this source too.
Original square sums of indistinguishable sources agree almost surely,
so their probability convergence transfers back.
\end{proof}
\begin{theorem}[Reverse Davis for general local martingales]
Under the usual conditions let $N$ be a zero-start strongly adapted
c\`adl\`ag local martingale. For any intrinsic quadratic variation and finite $T$,

\[
 E\sqrt{[N]_T}\leq7E N_T^*,
\] allowing infinite expectations.
\end{theorem}
\begin{proof}
At a common DDY stop $\rho_j$, the $L$ component is uniformly bounded almost
surely, so its original square sums converge in probability. Their difference
from the original $N^{\rho_j}$ sums converges pathwise to the FV correction.
Stopping consistency identifies the sum limit with $[N]_{T\wedge\rho_j}$.
Extract an almost-surely convergent subsequence and take roots. The grid
maxima are bounded by $(N^{\rho_j})_T^*$, so finite-grid reverse Davis and
Fatou give

\[
 E\sqrt{[N]_{T\wedge\rho_j}}\leq7E(N^{\rho_j})_T^*\leq7E N_T^*.
\]
Exhaustion makes the root eventually equal to $\sqrt{[N]_T}$ almost surely;
apply Fatou again and use uniqueness to cover any certificate.
\end{proof}
\begin{lemma}[Uniform integrability of convex square roots]
Fix a common stop and finite horizon. Let $E_i,B_i$ be the original and
$L$-component factorial square sums. For any forward weights $W$,
$\sqrt{W(E)_k}$ is uniformly integrable.
\end{lemma}
\begin{proof}
Let $A=\Var(Q^{\rho_j})_\infty$, which is integrable. The grid root bound gives
$\sqrt{E_i}\leq\sqrt{B_i}+A$, hence $E_i\leq2B_i+2A^2$.
Convex weights imply

\[
 W(E)_k\leq2W(B)_k+2A^2,\qquad
 \sqrt{W(E)_k}\leq2W(B)_k+2(1+A).
\]
The stopped $L$ is bounded, and its quadratic approximation estimates make
$B_i$ uniformly integrable, as are their forward convex averages.
The remaining majorant is one integrable function. Integrability of $A^2$
is not needed.
\end{proof}
\begin{lemma}[Terminal root convergence after closed stopping]
For a common stop $\rho$ and horizon $T$, let $V^L$ and its weights be from
the finite-horizon quadratic construction for $L^\rho$, and $E_k$ the same
weighted energies of $N^\rho$. With DDY correction $C$,

\[
 \sqrt{E_k(T)}\longrightarrow\sqrt{V^L_T+C_{T\wedge\rho}}
 \quad\text{in }L^1,
\]
and the limit is integrable.
\end{lemma}
\begin{proof}
Stopped grid increments are original increments evaluated at stopped times,
so the correction limit holds even when $t>\rho$.
The same weights give convergence of the $L$ energy and the correction,
hence almost-sure energy convergence. Continuity of square root and the
preceding uniform integrability permit Vitali's theorem.
The $L^2$ terminal value of $L^\rho$ supplies the finite-horizon quadratic data.
\end{proof}
\begin{lemma}[Identification with the original stopped quadratic variation]
If $V=[N]$, then
\[
 V^L_T+C_{T\wedge\rho}=V_{T\wedge\rho}
\]
 almost surely,
and the preceding $L^1$ limit is $\sqrt{V_{T\wedge\rho}}$.
\end{lemma}
\begin{proof}
Both $V^L$ and $[L]$ stopped at $\rho$ and $T$ are quadratic certificates
for $(L^\rho)^T$: their residuals are martingales and jumps are the source
squares. Uniqueness gives $V^L_T=[L]_{T\wedge\rho}$.
The indistinguishable identity $V=[L]+C$ also holds at the random time
$T\wedge\rho$. Substitute into the root limit.
\end{proof}
\begin{lemma}[Davis for a general local martingale]
Under the usual conditions let $N$ be a zero-start strongly adapted
c\`adl\`ag local martingale. For finite $T$,
\[
 E\sup_{t\leq T}|N_t|\leq6E\sqrt{[N]_T}
\]
 as extended
nonnegative expectations. A finite right side implies integrability of the maximum.
\end{lemma}
\begin{proof}
Fix a common stop $\rho$. A coarse factorial-grid maximum is below every
later refined-grid maximum. Average with the tail weights, use finite-grid
Davis, and then root concavity to bound its expectation by $6E\sqrt{E_k(T)}$.
The preceding $L^1$ convergence replaces this by
$6E\sqrt{[N]_{T\wedge\rho}}$ in the limit. Refine the coarse grid;
right continuity and monotone convergence give
\[
 E\sup_{t\leq T}|N^\rho_t|\leq6E\sqrt{[N]_{T\wedge\rho}}.
\]
The stopped maximum is integrable at this stage.
Now exhaust with $\rho_n$. Eventually the stopped maximum equals the original
one almost surely; the right side is at most $6E\sqrt{[N]_T}$ by monotonicity.
Fatou removes localization, and uniqueness removes the choice of certificate.
\end{proof}

\begin{lemma}[Infinite horizons and all-decomposition costs]
Put $R_N=\sup_{t\geq0}\sqrt{[N]_t}$. Then
\[
 E\sup_t|N_t|\leq6ER_N.
\]
For zero-start $Y$, consider every indistinguishable decomposition $Y=N+A$
with zero-start strongly adapted c\`adl\`ag local martingale $N$ and
strongly adapted c\`adl\`ag locally FV $A$. Define
\[
 \mathcal J(Y)=\inf_{Y=N+A}E[R_N+\Var(A)_\infty].
\]
Then $E\sup_t|Y_t|\leq6\mathcal J(Y)$, with no finite-cost assumption and
infimum of the empty family equal to infinity.
\end{lemma}
\begin{proof}
Increase integer horizons and use monotone convergence for the first bound.
The root supremum is measurable by integer-time detection and monotonicity.
For each decomposition, zero start gives

\[
 \sup_t|Y_t|\leq\sup_t|N_t|+\Var(A)_\infty.
\]
Take expectations, apply Davis, and take the infimum over all decompositions.
Existence and uniqueness of quadratic variation make its choices immaterial.
\end{proof}
\begin{lemma}[Closed stops and strict-prefix extension costs]
For any stopping time $\tau$, possibly infinite, $\mathcal J(Y^\tau)\leq\mathcal J(Y)$.
Set
\[
 \mathcal J_{\tau-}(Y)=\inf\{\mathcal J(L): L_t=Y_t\text{ for all }t<\tau\text{ a.s.}\}.
\]
Agreement uses one common null set. Then
\[
 \mathcal J_{\tau-}(Y)\leq\mathcal J(Y^\tau)\leq\mathcal J(Y),\qquad
 E\sup_{t<\tau}|Y_t|\leq6\mathcal J_{\tau-}(Y).
\]
The empty supremum at $\tau=0$ is zero.
\end{lemma}
\begin{proof}
Stop both components of each decomposition. Stopping consistency bounds the
root by $R_N$; variation does not increase. Take the infimum.
The stopped process itself extends the strict prefix. Every extension $L$
satisfies $\sup_{t<\tau}|Y_t|\leq\sup_t|L_t|$ almost surely; apply the
all-time estimate and take the extension infimum.
\end{proof}
\begin{lemma}[Triangle inequalities for both costs]
Under the same assumptions,
\[
 \mathcal J(X+Y)\leq\mathcal J(X)+\mathcal J(Y),\qquad
 \mathcal J_{\tau-}(X+Y)\leq\mathcal J_{\tau-}(X)+\mathcal J_{\tau-}(Y).
\]
The second inequality uses the same cutoff $\tau$, which need not be a
stopping time.
\end{lemma}
\begin{proof}
Add decompositions $X=N+A$ and $Y=M+B$ to obtain the admissible
decomposition $X+Y=(N+M)+(A+B)$.
The quadratic-root triangle inequality gives $R_{N+M}\leq R_N+R_M$
almost surely, while variation subadditivity gives
$\Var(A+B)_\infty\leq\Var(A)_\infty+\Var(B)_\infty$.
All-time variation is the supremum of integer-horizon variations and is
measurable for right-continuous adapted paths. We may therefore add the
expectations, bounding the sum decomposition's cost by the two original
costs. Taking independent infima over the two decomposition families proves
the first inequality.
Adding two global extensions of the same strict prefix gives an extension
of the sum target. Apply the first inequality and take independent infima
over extensions for the second. Infinite costs and empty families are included.
\end{proof}
\begin{lemma}[Negation and successive differences]
Under the same assumptions, $\mathcal J(-X)=\mathcal J(X)$ and
$\mathcal J_{\tau-}(-X)=\mathcal J_{\tau-}(X)$.
For one fixed $\tau$, therefore,
\[
 \mathcal J_{\tau-}(Z_k-X)\longrightarrow0
 \quad\Longrightarrow\quad
 \mathcal J_{\tau-}(Z_{k+1}-Z_k)\longrightarrow0.
\]
\end{lemma}
\begin{proof}
A quadratic certificate for $N$ also satisfies the conditions for $-N$:
the square residual is unchanged and squared jumps ignore the sign.
Variation is invariant under negation too, so $(N,A)\mapsto(-N,-A)$
preserves costs. Taking infima and applying negation twice proves the
identity for $\mathcal J$. Negating extensions of the same prefix gives
the corresponding identity for the outer infimum.
The triangle inequality and symmetry give
\[
 \mathcal J_{\tau-}(Z_{k+1}-Z_k)
 \leq\mathcal J_{\tau-}(Z_{k+1}-X)+\mathcal J_{\tau-}(Z_k-X).
\]
Combine this with nonnegativity for the limit. Infinite costs are allowed,
but the lemma does not itself construct a common stop.
\end{proof}
\begin{lemma}[Components from finite strict-prefix cost]
If $\tau\leq T$ is a finite stop and $\mathcal J_{\tau-}(X)<r$, there are
$X'\sim X$, a local martingale $N$, and globally BV $A$, both adapted and
c\`adl\`ag, such that $X'_t=N_t+A_t$ pointwise for $t<\tau$, and
\[
 E\sup_{t\leq T}|N_{t\wedge\tau}|+E\Var_{[0,T]}(A^{\tau-})<6r.
\]
\end{lemma}
\begin{proof}
Choose an extension $Y$ and decomposition $Y=N+B$ with cost less than $r$.
Put $A=B^T$. Its strict-prefix variation is bounded by the original total
variation, and all-time Davis bounds the stopped martingale maximum.
Define $X'$ as $N+A$ on the strict prefix and as $X$ elsewhere.
On the common event of extension and decomposition agreement, $X'=X$ at
all times. This gives both the exact prefix identity and indistinguishability.
\end{proof}
\begin{lemma}[Component sequences from summable costs]
Under the usual conditions let $\tau\leq T$ be one finite stopping time.
If $\sum_k\mathcal J_{\tau-}(X_k)<\infty$, one can choose the preceding
exact representatives $X'_k$ and components $(N_k,A_k)$ with
\[
 \sum_k\left(E\sup_{t\leq T}|N_{k,t\wedge\tau}|+
 E\Var_{[0,T]}(A_k^{\tau-})\right)
 \leq6\left(\sum_k\mathcal J_{\tau-}(X_k)+1\right)<\infty.
\]
On one full-measure set $X'_{k,t}=X_{k,t}$ for every $k,t$, and
\[
 \sum_k\left(\sup_{t\leq T}|N_{k,t\wedge\tau}|+
 \Var_{[0,T]}(A_k^{\tau-})\right)<\infty.
\]
\end{lemma}
\begin{proof}
Put $c_k=\mathcal J_{\tau-}(X_k)$ and $\varepsilon_k=2^{-k-1}$.
Each $c_k$ is finite, so $c_k<c_k+\varepsilon_k$. The preceding lemma
supplies representatives and components with cost less than
$6(c_k+\varepsilon_k)$.
Sum and use $\sum_k\varepsilon_k=1$ to obtain the expectation bound.
The martingale maxima and strict-prefix variations are measurable;
Tonelli gives pathwise finiteness of their sum almost surely.
Intersect with the countably many indistinguishability events to retain
agreement with all original targets, at every time, on that same set.
\end{proof}
The common stop and summability are inputs here. These arbitrary decomposition
components still require a separate bridge to actual original-price integrals.
\begin{lemma}[Strict-prefix comparison and boundary jumps]
For finite $\tau\leq T$,

\[
 \mathcal J(X^{\tau-})\leq2\mathcal J(X).
\]
If $X_0=0$, then

\[
 \mathcal J_{\tau-}(X)\leq\mathcal J(X^{\tau-})\leq2\mathcal J_{\tau-}(X).
\]
For zero-start adapted c\`adl\`ag $X$,
\[
 \mathcal J(X^\tau)\leq2\mathcal J_{\tau-}(X)+E|\Delta X_\tau|.
\]
Infinite costs are allowed.
\end{lemma}
\begin{proof}
For $X=N+A$, stopping $A$ at $T$ leaves the prefix unchanged.
Use the decomposition
$X^{\tau-}=N^\tau+(A^{\tau-}-B^\tau(N))$, where $B^\tau$ is the boundary
jump process. The jump-square identity and monotonicity give
$|\Delta N_\tau|\leq\sqrt{[N]_\tau}$.
The stopped root does not increase; variation increases by at most this
jump, so the cost is at most twice the original. Take infima.
An extension agreeing before $\tau$ has the same left limit there; at
$\tau=0$ use zero start. Apply the same estimate to extension decompositions
and take infima. The reverse bound follows because the strict prefix is
itself an extension. Finally $B^\tau(X)$ is adapted zero-start FV with
cost at most $E|\Delta X_\tau|$. Add it to $X^{\tau-}$ and use the triangle
inequality.
\end{proof}
These results concern the zero-start core and its extension costs. The next
subsection constructs common prelocal control and supplies the algebraic
inputs for the main sequence. Nonzero initial values are separate coordinates.
A finite $M^2$ quadratic root is not itself exact $j^1$, and the stopping-jump
seminorm is distinct from the root-plus-variation cost. Quadratic variation
does not detect constant martingales, which is why components are centered.
